% Pursuing Stacks (cotes 134-2 to 134-8): archive cote 134-4, site page 37 (PDF page 38), typescript page 125.
% typescript. Draft compared: ps3.tex:1841-1908.
% AI-transcribed by Claude Opus 5.5 (claude-opus-5-5, adaptive thinking, effort high), Message Batches API, batch ps2 (msgbatch_01NVwr1Yuo8qJf5rykeGBXPK), 2026-11-12; not yet checked by eye.
\scanpage{134-4/37}
\mspage{125}
objects. O\unsure{ne} way is to take for $M'$ any quotient category of $M$, by an
equivalence relation which is the discrete one on objects - thus $M'$ cor-%
responds just to an equivalence relation $R$ in $F\hand{\ell}(M)$, compatible with
\xout{\unsure{compositi}} the source and target maps and with compositions. T\unsure{h}e other is
to take as $M'$ a localization with respect to some $W\hand{\subset}F\hand{\ell}(M)$, and if we
take $W$ strongly saturated we get a bijective correspondance between those
$M'$ and the \xout{corresponding} set of strongly saturated subsets $W$ of $F\hand{\ell}(M)$.

It would'nt be any more reasonable to call an arbitrary $R$ as above,
corresponding to an arbitrary quotient category $M'$ with same objects as
$M$, a ``homotopy relation'' o\over{\unsure{f}}{n} $M$ (as I did though yesterday), as it would
be to call an arbitrary strongly saturated $W\hand{\subset}F\hand{\ell}(M)$ a ``homotopism struc-%
ture'' on $M$ (as I nearly did yesterday, but then rectified in the stride)\hand{.}
\uline{The characteristic flavor of homotopy theory comes in}, \uline{when we got an}
$M'$ \uline{which is} \uline{both} \uline{a quotient category and a localization of} $M$. Thus \xout{neithe}
neither approach, via $R$ or via $W$, is any more contained in the other\hand{,}
than \xout{to} the converse. \over{\unsure{w}}{W}e should regard the \uline{homotopy structure} on $M$ to
be embodied in the basic functor
\[ M \longrightarrow M' \quad , \]
which is a description where no choice yet is made between the two possib-%
le descriptions of $M'$, either by an $R$, or by a $W$. If we describe $M'$ in
terms of $R$, the extra assumption to make on $R$ for calling it a ``\uline{homotopy}
\uline{relation}'', is \xout{the\unsure{xxxx}} that the \inserted[typed]{canonical} functor
\[ M \longrightarrow M/R \]
should be localizing. Alternatively, describing $M'$ by $W$, the extra
assumption to make on $W$ (as we did yesterday) in order to call $W$ a \uline{homoto-}%
\uline{pism} \uline{structure} on $M$, is that the canonical \inserted[typed]{localization} functor
\[ M \longrightarrow W^{-1}M \]
be essentially a passage-to-quotient functor, namely \inserted[typed]{sur}\xout{bi}jective on arrows
(as we know already it is bijective on objects). Thus the set of all
homotopy relations on $M$ is in one to one correspondance with the set of
all homotopism structures on $M$, and if we denote these sets (in accordan-%
ce with yesterday's provisional notations) $\mathrm{Hom}_1(M)$ and $\mathrm{Hom}_2(M)$, we get
thus a \uline{bijective} correspondance
\[ \mathrm{Hom}_1(M) \longleftrightarrow \mathrm{Hom}_2(M) \qquad . \]
The set of homotopy structures $\mathrm{Hom}(M)$ on $M$ may be either defined as the
usual quotient set defined by the previous two-member system of transitive
bijections between sets, or more substantially, as a set of isomorphism
classes of \unsure{\xout{``}}categories $M'$ ``under $M$'' i.e. endowed with a functor $M\to M'$,
\inserted{and} subject to the following two extra conditions:

\margin{37}
